\documentclass[../../main.tex]{subfiles}

\begin{document}

\section*{附加习题：拓扑群}\label{section:ch02-附加习题拓扑群}

\addcontentsline{toc}{section}{附加习题：拓扑群}

在这些习题中，我们研究拓扑群及其某些性质。商拓扑的名称便是由构造拓扑群关于其子群的商群而来。

拓扑群（topological group）$G$ 既是一个群又是一个满足 $T_{1}$ 公理的拓扑空间，并且将 $x \times y$ 映为 $x \cdot y$ 的从 $G \times G$ 到 $G$ 的映射以及将 $x$ 映为 $x^{-1}$ 的从 $G$ 到 $G$ 的映射都是连续的。在下面的题目中，用 $G$ 表示拓扑群。

\begin{exercise}\label{exercise:ch02-附加习题-1}
设 $H$ 是一个群，同时也是一个满足 $T_{1}$ 公理的拓扑空间。证明：$H$ 是一个拓扑群当且仅当将 $x \times y$ 映成 $x \cdot y^{-1}$ 的从 $H \times H$ 到 $H$ 的映射是连续的。
\end{exercise}

\begin{exercise}\label{exercise:ch02-附加习题-2}
证明以下都是拓扑群：

(a) $(\mathbb{Z}, +)$。

(b) $(\mathbb{R}, +)$。

(c) $(\mathbb{R}_{+}, \cdot)$。

(d) $(S^{1}, \cdot)$，其中 $S^{1}$ 是使得 $|z| = 1$ 的所有复数 $z$ 组成的空间。

(e) 在矩阵乘法运算下的一般线性群 $GL(n)$。（$GL(n)$ 是所有非蜕化 $n \times n$ 阶矩阵的集合，其拓扑是作为 $n^{2}$ 维欧氏空间的子集按显然的方式给出的。）
\end{exercise}

\begin{exercise}\label{exercise:ch02-附加习题-3}
设 $H$ 是 $G$ 的一个子空间。证明：如果 $H$ 又是 $G$ 的一个子群，那么 $H$ 和 $\overline{H}$ 两者都是拓扑群。
\end{exercise}

\begin{exercise}\label{exercise:ch02-附加习题-4}
设 $\alpha$ 是 $G$ 的一个元素。证明：用
\begin{align*}
f_{\alpha}(x) = \alpha \cdot x \quad \text{以及} \quad g_{\alpha}(x) = x \cdot \alpha
\end{align*}
定义的映射 $f_{\alpha}, g_{\alpha}: G \rightarrow G$ 是 $G$ 的一个同胚。证明 $G$ 是一个齐性空间。（即对于 $G$ 的任意两个点 $x$ 和 $y$，存在一个从 $G$ 到 $G$ 上的同胚将 $x$ 映成 $y$。）
\end{exercise}

\begin{exercise}\label{exercise:ch02-附加习题-5}
设 $H$ 是 $G$ 的一个子群。若 $x \in G$，定义 $xH = \{x \cdot h \mid h \in H\}$，称之为 $H$ 在 $G$ 中的一个左陪集（left coset）。设 $G/H$ 表示 $G$ 中 $H$ 的左陪集族，它是 $G$ 的一个分拆。在 $G/H$ 上赋予商拓扑。

(a) 证明：若 $\alpha \in G$，前一题中的映射 $f_{\alpha}$ 诱导出 $G/H$ 的一个同胚，它将 $xH$ 映为 $(\alpha \cdot x)H$。证明 $G/H$ 是一个齐性空间。

(b) 证明：若 $H$ 关于 $G$ 的拓扑是一个闭集，则单点集在 $G/H$ 中都是闭的。

(c) 证明：商映射 $p: G \rightarrow G/H$ 是开的。

(d) 证明：若 $H$ 关于 $G$ 的拓扑是一个闭集，并且是 $G$ 的一个正规子群，则 $G/H$ 是拓扑群。
\end{exercise}

\begin{exercise}\label{exercise:ch02-附加习题-6}
整数集 $\mathbb{Z}$ 是 $(\mathbb{R}, +)$ 的正规子群。商 $\mathbb{R}/\mathbb{Z}$ 是一个你所熟悉的拓扑群，它是哪个拓扑群？
\end{exercise}

\begin{exercise}\label{exercise:ch02-附加习题-7}
若 $A$ 和 $B$ 都是 $G$ 的子集，用 $A \cdot B$ 表示所有点 $a \cdot b$ 的集合，其中 $a \in A$，$b \in B$。用 $A^{-1}$ 表示所有点 $a^{-1}$ 的集合，其中 $a \in A$。

(a) 单位元 $e$ 的一个邻域 $V$ 称为对称的（symmetric），若 $V = V^{-1}$。若 $U$ 为 $e$ 的一个邻域，证明：存在一个 $e$ 的对称邻域 $V$，使得 $V \cdot V \subset U$。[提示：若 $W$ 为 $e$ 的一个邻域，则 $W \cdot W^{-1}$ 便是对称的。]

(b) 证明 $G$ 是一个 Hausdorff 空间。事实上，只要证明：若 $x \neq y$，则存在 $e$ 的邻域 $V$，使得 $V \cdot x$ 与 $V \cdot y$ 无交。

(c) 证明：$G$ 满足以下称为正则性公理（regularity axiom）的分离公理：对于任意给定的闭子集 $A$ 和 $A$ 外的一个点 $x$，存在分别包含 $A$ 和 $x$ 的无交的两个开集。[提示：存在 $e$ 的一个邻域 $V$，使得 $V \cdot x$ 与 $V \cdot A$ 无交。]

(d) 设 $H$ 是 $G$ 的一个子群并且关于 $G$ 的拓扑是闭的。$p: G \rightarrow G/H$ 为商映射。证明 $G/H$ 满足正则性公理。[提示：当 $A$ 为饱和集时，考察 (c) 小题的证明。]
\end{exercise}

\end{document}